\subsection{李群的自同构群（超度量情形）}

定理 2. 当 K 是超度量且局部紧的、G 是紧李群时，定理 1 的断言 (i)、(ii)、(iii)、(v) 和 (vi) 成立。

\begin{enumerate}
\def\labelenumi{(\alph{enumi})}
\item
  (i) 中所考虑的解析结构的唯一性是显然的。
\item
  假设 G 是由可赋范李代数 L 定义的李群。则 G 是 L 中的一个开闭球。令 \(w \in \operatorname{Aut} G\)。于是 \(\operatorname{L}(w)\) 在 0 的一个邻域内与 w 重合。令 \(x \in G\)。令 p 为剩余域的特征。于是 \(p^{n}x\) 当 n 趋于 \(+\infty\) 时趋于 0。因此存在 n，使得 \(w(p^{n}x) = \operatorname{L}(w)(p^{n}x)\)。所以
\end{enumerate}

\[
\begin{array}{r l} p ^ {n} w (x) & = w (x) ^ {p ^ {n}} = w (x ^ {p ^ {n}}) = w (p ^ {n} x) \\ & = \mathrm{L} (w) (p ^ {n} x) = p ^ {n} \mathrm{L} (w) (x) \end{array}
\]

从而 \(w(x) = \mathrm{L}(w)(x)\)。因此，\(w = \mathrm{L}(w)|\mathrm{G}\)。

令 \(\Gamma\) 为这样的 \(\gamma \in \operatorname{Aut} L(G)\) 的集合，满足 \(\gamma(G) = G\)。由于 \(G\) 在 \(L(G)\) 中既开又紧，\(\Gamma\) 是 \(\operatorname{Aut} L(G)\) 的开子群。由上面所述，Aut \(G\) 可与 \(\Gamma\) 识别，从而 Aut \(G\) 上存在李群结构，在此结构下定理 1 的性质 (i)、(ii)、(iii) 和 (vi) 显然成立。性质 (v) 由第 1 号命题 1 和 3 得出。

\begin{enumerate}
\def\labelenumi{(\alph{enumi})}
\setcounter{enumi}{2}
\tightlist
\item
  转到一般情形。由第 7 节第 1 号命题 1，存在一个开紧子群 \(G_0\)，它是 \(G\) 的子群，属于 (b) 中所考虑的类型。于是 \(G\) 由 \(G_0\) 和有限个元素 \(x_1, x_2, \ldots, x_n\) 生成。令 \(\text{Aut}_1G\) 为这样的 \(u \in \text{Aut}G\) 的集合：满足 \(u(G_0) = G_0\) 且 \(u(x_iG_0) = x_iG_0\)（\(1 \leqslant i \leqslant n\)）。按照定理 1 的证明 (c) 部分，定义 \(P\)，即 \(\text{Aut}G_0\) 与 \(G_0^n\) 的半直积，以及单同态 \(\zeta\)，它从 \(\text{Aut}_1G\) 到 \(P\)，其像在 \(P\) 中是闭的。
\end{enumerate}

(d)、(e)：论证完全同定理 1 证明的 (d)、(e) 部分，只需将 \(\mathbf{R}\) 换成 \(\mathbf{Q}_p\)，并用命题 3 代替命题 2。

注。若 \(\mathbf{K} = \mathbf{Q}_p\) 且李群 \(\mathbf{G}\) 由一个紧子集生成（参见~练习 2），则定理 1 的断言 (i)、(ii)、(iii) 和 (vi) 仍然成立，但 (v) 不成立（练习 3）。

